\documentclass[11pt]{article}
\usepackage[T1]{fontenc}
\usepackage[utf8]{inputenc}
\usepackage[margin=1in]{geometry}
\usepackage{amsmath,amssymb,amsthm}
\usepackage{microtype}
\usepackage{setspace}
% 9-inch text block / 21.76pt baseline: about 30 prose lines per page.
% Extra room accommodates headings and displayed mathematics.
\setstretch{1.6}
\raggedbottom
\usepackage[hidelinks]{hyperref}
\newcommand{\E}{\mathbb E}
\newcommand{\Pb}{\mathbb P}
\DeclareMathOperator{\Beta}{Beta}
\DeclareMathOperator{\Var}{Var}
\newcommand{\Bfun}{\mathrm{B}}
\newtheorem{proposition}{Proposition}
\newtheorem{lemma}{Lemma}
\newtheorem{corollary}{Corollary}
\newtheorem{conjecture}{Conjecture}
\theoremstyle{definition}
\newtheorem{definition}{Definition}
\newtheorem{remark}{Remark}
\allowdisplaybreaks[1]
\title{Two Strategic Sellers under Forgetting:\\
The Fluid Model}
\author{Standalone theoretical note}
\date{October 8, 2026}

\begin{document}
\maketitle

Two sellers choose hidden product quality while a customer learns from
sales outcomes and gradually forgets them. We study open-loop Nash
equilibrium in a deterministic fluid approximation. The main result is an
exact test for symmetric stationary equilibrium, conditional on
log-concavity of marginal demand: reputation must be a fixed point of the
frozen-payoff peak map, and holding must beat using the cheap product
forever.

Sections~\ref{sec:model}--\ref{sec:transformation} set up the game and
transform its payoff. Section~\ref{sec:frozen-single-interval} gives the
single-interval structure used in the stationary test of
Section~\ref{sec:fluid-steady}. Section~\ref{sec:moving} treats candidate
moving paths and their verification. All proofs and derivations are in the
appendices. We prove the demand-shape conjecture for rivals with integer
Beta shape parameters, real shapes in $[1,2]^2$, and nearby parameters.
Its general form, equilibrium
existence, and convergence remain open. Numerical checks in
Appendix~\ref{app:demand-numerics} are exploratory, not proofs.

\section{Model}\label{sec:model}

\paragraph{Primitives.}
Two symmetric sellers $j\in\{A,B\}$ each offer products $i\in\{1,2\}$. A
product-$i$ sale succeeds with probability $p_i$ and costs the seller $c_i$,
where $0<p_1<p_2<1$ and $0\leq c_1<c_2<R$, and $R$ is the revenue per sale.
Write $\Delta p=p_2-p_1$ and $\Delta c=c_2-c_1$. The sellers discount by
$\gamma\in(0,1)$ per period. The product used is never observed by the
customer.

The sellers observe both sellers' outcomes, but not the rival's product.
The customer's information, and hence the reputation state, is public.

\paragraph{Learning and forgetting.}
For each seller, the customer maintains effective success and failure
counts $(S,F)$, initially $(0,0)$, with posterior $\Beta(1+S,1+F)$.
An observation updates the chosen seller's counts to $(\rho S+1,\rho F)$
after success and $(\rho S,\rho F+1)$ after failure; the other seller's
counts do not change. We assume $\rho\in(1/2,1)$ and set $n=(1-\rho)^{-1}$.

After $k$ observations, the normalized total is
$(1-\rho)(S+F)=1-\rho^k$. On its limiting unit segment, reputation is a
success share $x\in[0,1]$, with posterior $\Beta(1+nx,1+n(1-x))$ and updates
\begin{equation}\label{eq:share-maps}
f_1(x)=\rho x+1-\rho,\qquad f_0(x)=\rho x.
\end{equation}
We work on the limiting state space $(x,y)\in[0,1]^2$, with $x$ for seller
$A$ and $y$ for seller $B$. Appendix~\ref{app:model} justifies this state
space and records the demand properties used below.

\paragraph{Customer choice and demand.}
In each period the customer draws independently
$\widetilde\theta_A\sim\Beta(1+nx,1+n(1-x))$ and
$\widetilde\theta_B\sim\Beta(1+ny,1+n(1-y))$ (Thompson sampling), and
chooses the seller with the larger draw. Seller $A$'s demand is
\begin{equation}\label{eq:boundary-demand}
D_A(x,y)=\Pb\bigl(\widetilde\theta_A>\widetilde\theta_B\bigr),
\qquad D_B(x,y)=1-D_A(x,y).
\end{equation}
Some seller is chosen in every period. By symmetry,
$D_A(x,y)=1-D_A(y,x)$ and $D_A(x,x)=1/2$. Since both posteriors have full
support, $D_A$ is smooth on $[0,1]^2$ with values in $(0,1)$.

Demand $D_A$ increases with $x$ and decreases with $y$.

\paragraph{Single-seller benchmark.}
For comparison we use the model in which seller $A$ faces an outside option
of known quality $p_0\in(0,1)$, so that demand is
\begin{equation*}
S_{p_0}(x)=\Pb\bigl(\Beta(1+nx,1+n(1-x))>p_0\bigr).
\end{equation*}
This is the case of a rival whose posterior is a point mass at $p_0$.

\section{The fluid model}\label{sec:fluid-derivation}

Let $\lambda=-\log\rho$, $\delta=-\log\gamma$, and let
$u_j(t)\in[0,1]$ be seller $j$'s fraction of product~2. The fluid
approximation replaces random reputation increments by their drift:
\begin{equation}\label{eq:duopoly-fluid}
\begin{aligned}
\dot x&=\lambda D_A(x,y)\bigl(p_1+u_A\Delta p-x\bigr),\\
\dot y&=\lambda D_B(x,y)\bigl(p_1+u_B\Delta p-y\bigr),
\end{aligned}
\end{equation}

Seller $A$ maximizes
\begin{equation}\label{eq:payoff}
J_A(u_A,u_B;x_0,y_0)=\int_0^\infty e^{-\delta t}
D_A(x(t),y(t))(R-c_1-u_A(t)\Delta c)\,dt,
\end{equation}
with the analogous objective for $B$. Controls are measurable time paths.
Each pair induces a unique absolutely continuous trajectory that remains
in $[0,1]^2$.

The approximation is most natural when $\rho$ is close to one. Its formal
embedding scales period payoffs by $dt$, discounting by $\gamma^{dt}$, and
reputation updates by $\rho^{dt}$, while keeping demand's original $n$.
It does not model the initial buildup of the effective counts.
Appendix~\ref{app:fluid} gives the derivation.

\section{Solution concepts}\label{sec:fluid-solution-concepts}

Open-loop play commits each seller to a time path from a specified initial
state, as in the ex ante sequence of choices in \cite{hu2026}.

\begin{definition}[Open-loop Nash equilibrium]
A pair of control paths $(u_A^*(\cdot),u_B^*(\cdot))$ is an open-loop Nash
equilibrium from $(x_0,y_0)$ if $u_A^*$ maximizes
$J_A(\cdot,u_B^*;x_0,y_0)$ and $u_B^*$ maximizes
$J_B(u_A^*,\cdot\,;x_0,y_0)$.
\end{definition}

Fixing $u_B(\cdot)$ does not fix $y(\cdot)$: a deviation by $A$ changes
$B$'s demand and thus its updating speed. The best-response problem has
state $(x,y)$ even when the proposed equilibrium is symmetric. A unilateral
deviation generally leaves the diagonal.

\paragraph{The frozen-rival problem.}
When the rival holds at $y\in[p_1,p_2]$ by committing to
$u_B=(y-p_1)/\Delta p$, its reputation remains $y$ under every deviation.
Seller $A$ then faces the one-dimensional problem
\begin{equation}\label{eq:frozen-rival}
\sup_{u(\cdot)}\int_0^\infty e^{-\delta t}D_A\bigl(x(t),y\bigr)
\bigl(R-c_1-u(t)\Delta c\bigr)\,dt
\quad\text{s.t.}\quad
\dot x=\lambda D_A(x,y)\bigl(p_1+u\Delta p-x\bigr).
\end{equation}

This is the best-response problem for stationary open-loop play. For other
fixed $y\in[0,1]$, it is an auxiliary problem of the same form as the
single-seller benchmark.

\paragraph{Feedback comparison.}
A feedback Nash equilibrium instead specifies state-dependent controls
that are mutual best responses from every state. It is the fluid analogue
of Markov perfect equilibrium. Deviations can change the rival's action,
so the stationary test below does not generally apply. The feedback HJB
and the additional strategic term are recorded in
Appendix~\ref{app:feedback}.

Neither equilibrium's existence is automatic. In the open-loop problem,
nonlinear state dynamics can make payoffs nonconcave in one's control and
best-response sets nonconvex.

\section{A payoff transformation}\label{sec:transformation}

Define
\begin{equation*}
q=p_1+\frac{(R-c_1)\Delta p}{\Delta c}>p_2,
\qquad
F_A(x,y)=\frac{\Delta c}{\Delta p}D_A(x,y)(q-x)
-\frac{\delta\Delta c}{\lambda\Delta p}x.
\end{equation*}
The discounted payoff has the equivalent representation
\begin{equation}\label{eq:duopoly-transform}
J_A(u_A,u_B;x_0,y_0)=\frac{\Delta c}{\lambda\Delta p}x_0
+\int_0^\infty e^{-\delta t}F_A(x(t),y(t))\,dt.
\end{equation}
The initial-state term is independent of the controls. Product choice
therefore only steers the state in the transformed problem: the direct
quality cost is absorbed by valuing each unit of reputation drift at
$\Delta c/(\lambda\Delta p)$. With a holding rival, this leaves a
one-dimensional payoff $F_A(\cdot,y)$ whose shape determines the best
response. Appendix~\ref{app:transform} proves the identity.

\section{Single-interval structure in the frozen-rival problem}
\label{sec:frozen-single-interval}

Fix the rival's reputation $y$ and write $D(x)=D_A(x,y)$ and
$F(x)=F_A(x,y)$. The key condition concerns marginal demand.

\begin{proposition}[Single interval under log-concave marginal demand]
\label{prop:frozen-reduction}
If $x\mapsto\partial_xD_A(x,y)$ is log-concave on $[0,1]$, then the
frozen-rival problem \eqref{eq:frozen-rival} admits an optimal stationary
policy that uses product~2 on at most one interval of $x$. The nonconstant
patterns, in increasing order of $x$, are $1\to2$, $2\to1$, and
$1\to2\to1$, with holding mixtures at attractive thresholds in
$(p_1,p_2)$.
\end{proposition}

Under this hypothesis, $F$ is either nonincreasing or decreases, then
increases, then decreases, with absent segments allowed. In particular,
there is at most one interior local maximum. This shape reduces the
stationary equilibrium test to a comparison between holding at the peak
and descending toward $p_1$. The proof and the underlying control lemma
are in Appendix~\ref{app:control}.

\begin{corollary}[Single seller]\label{cor:single-seller}
In the single-seller benchmark, the fluid problem admits an optimal
stationary policy that uses product~2 on at most one interval of $x$.
\end{corollary}

The demand property can also be proved exactly for a class of uncertain
rivals. The next result allows the sellers' posterior totals to differ.

\begin{proposition}[Integer and real Beta shape parameters]\label{prop:integer-rival}
Let $N>0$ and $\alpha,\beta\geq1$, and define
\begin{equation*}
\mathcal D(a)=\Pb\bigl(\Beta(a,N-a)>\Beta(\alpha,\beta)\bigr),
\qquad 0<a<N,
\end{equation*}
with independent draws. Suppose either that both rival shapes are integers
or that $1\leq\alpha,\beta\leq2$, with no integrality requirement in the
second case. Then $\mathcal D'(a)>0$ and $a\mapsto\mathcal D'(a)$ is strictly
log-concave on $(0,N)$ if $\alpha+\beta>2$. For the uniform rival,
$\alpha=\beta=1$, it equals $1/N$. In the real-shape range $[1,2]^2$,
$\mathcal D'$ is also concave.
\end{proposition}

In the common-$n$ model, this proves the conjectured shape for every
integer $n=m\geq1$ and $y=k/m$, $k=0,\ldots,m$, uniformly over all
$x\in[0,1]$. The economic restriction $\rho>1/2$ selects $m\geq3$.
The real-shape case additionally covers $ny\leq1$ and $n(1-y)\leq1$:
every $y$ if $0<n\leq1$, and $1-1/n\leq y\leq1/n$ if $1<n\leq2$.
Although this explicit range lies outside the economic restriction $n>2$,
strictness extends every proved parameter point to a neighborhood, as
described in Corollary~\ref{cor:integer-neighborhood}. The stationary test
is unconditional on demand shape wherever these results apply.
Appendix~\ref{app:integer-rival} gives the polynomial proof for integer
shapes and a Dirichlet-coupling proof for real shapes in $[1,2]^2$.
Arbitrary real shapes above two remain open.

\begin{conjecture}[Log-concave marginal demand in the duopoly]
\label{conj:duopoly-logconcave}
For every $n>0$ and $y\in[0,1]$, $x\mapsto\partial_xD_A(x,y)$ is
log-concave on $[0,1]$.
\end{conjecture}

Appendix~\ref{app:demand} collects further supporting results. The condition
also holds against a sufficiently concentrated rival
(Proposition~\ref{prop:concentrated-rival}). A finite-observation analogue
is also log-concave (Proposition~\ref{prop:discrete-jump-logconcave}), but
this does not establish the infinitesimal claim. An exact representation
reduces the conjecture to the variance inequality
\eqref{eq:logconcavity-variance}; a large-$n$ approximation provides further
heuristic support.

The single-interval result concerns reputation in a frozen-rival problem.
It does not bound the number of switches along a moving equilibrium path:
applying the frozen policy at each current $y(t)$ need not be a best response.

\section{Stationary open-loop Nash equilibrium}\label{sec:fluid-steady}

A stationary profile starts at a rest point $(x^*,y^*)\in[p_1,p_2]^2$
and uses constant holding controls
$u_A=(x^*-p_1)/\Delta p$ and $u_B=(y^*-p_1)/\Delta p$.
It is an open-loop equilibrium exactly when holding solves each seller's
frozen-rival problem. This is an equilibrium from the rest point, not a
claim about play from other initial states.

Under Section~\ref{sec:frozen-single-interval}'s shape hypothesis, let
$m(y)$ denote the interior peak of $F_A(\cdot,y)$ when it exists.
For a proposed symmetric reputation $z$, define
\begin{equation*}
\dot X_1(t;z)=\lambda D_A(X_1(t;z),z)(p_1-X_1(t;z)),
\qquad X_1(0;z)=z,
\end{equation*}
and the transformed value of using product~1 forever,
\begin{equation}\label{eq:stationary-low-value}
L_z(z)=\int_0^\infty e^{-\delta t}F_A(X_1(t;z),z)\,dt.
\end{equation}

\begin{proposition}[Symmetric interior stationary open-loop equilibrium]
\label{prop:stationary-open-loop}
Fix $z\in(p_1,p_2)$ and suppose $x\mapsto\partial_xD_A(x,z)$ is
log-concave on $[0,1]$. The constant control paths
\begin{equation*}
u_A^*(t)=u_B^*(t)=\frac{z-p_1}{\Delta p},\qquad t\geq0,
\end{equation*}
form an open-loop Nash equilibrium from $(z,z)$ if and only if
\begin{equation}\label{eq:stationary-exact-test}
z=m(z)
\qquad\text{and}\qquad
\frac{F_A(z,z)}{\delta}\geq L_z(z).
\end{equation}
Equivalently, the peak condition is
\begin{equation}\label{eq:stationary-peak-test}
(q-z)\partial_xD_A(z,z)=\frac12+\frac{\delta}{\lambda},
\qquad \partial_{xx}F_A(z,z)<0.
\end{equation}
The derivative in the second condition holds the rival's reputation fixed.
\end{proposition}

The proof is in Appendix~\ref{app:stationary}. To apply the test, solve
\eqref{eq:stationary-peak-test} for interior peaks and retain those that
pass \eqref{eq:stationary-exact-test}. The first-order equality alone is
insufficient: it can select a valley or a peak where abandoning quality
is preferable. The result is conditional on the demand-shape hypothesis
at the candidate and does not guarantee existence or uniqueness.

\paragraph{Boundary and asymmetric profiles.}
Let $W^{\mathrm{fr}}(s;r)$ be the optimal transformed value from own
reputation $s$ against a rival fixed at $r$. A holding profile at
$(x^*,y^*)$ is an equilibrium exactly when
\begin{equation}\label{eq:general-stationary-test}
W^{\mathrm{fr}}(x^*;y^*)=\frac{F_A(x^*,y^*)}{\delta},
\qquad
W^{\mathrm{fr}}(y^*;x^*)=\frac{F_A(y^*,x^*)}{\delta}.
\end{equation}
This reduction needs no log-concavity. Under the shape hypothesis, each
interior seller must be at its frozen-payoff peak and pass its product-1
value comparison. The symmetric boundary candidates $(p_1,p_1)$ and
$(p_2,p_2)$ instead require one-sided deviation tests. The value construction
in Appendix~\ref{app:control} applies to these cases as well.

\section{Moving open-loop paths}\label{sec:moving}

Fix the rival's path $u_B^*(t)$ and abbreviate
$D=D_A(x,y)$, $h_A=p_1+u_A\Delta p-x$, and
$h_B=p_1+u_B^*(t)\Delta p-y$. Let $\xi$ and $\eta$ be the
current-value costates of own and rival
reputation in the transformed problem. Its Hamiltonian is
\begin{equation*}
\mathcal H_A=F_A+\lambda\xi D h_A+\lambda\eta(1-D)h_B.
\end{equation*}
Under the regularity and normality conditions for the infinite-horizon
maximum principle, the adjoint equations are
\begin{equation}\label{eq:transformed-adjoints}
\begin{aligned}
\dot\xi&=(\delta+\lambda D)\xi-\partial_xF_A
 -\lambda\partial_xD\,(\xi h_A-\eta h_B),\\
\dot\eta&=(\delta+\lambda(1-D))\eta-\partial_yF_A
 -\lambda\partial_yD\,(\xi h_A-\eta h_B).
\end{aligned}
\end{equation}
Hamiltonian maximization gives $u_A=1$ when $\xi>0$ and $u_A=0$ when
$\xi<0$. On a singular interval, $\xi=\dot\xi=0$, hence
\begin{equation}\label{eq:open-loop-singular}
\partial_xF_A=\lambda\eta\,\partial_xD_A
\bigl(p_1+u_B^*(t)\Delta p-y\bigr).
\end{equation}
This condition reduces to $\partial_xF_A=0$ when the rival holds.
Otherwise the rival's drift and its shadow value enter, and singularity
does not imply $\dot x=0$. Determining a singular control requires further
conditions that keep \eqref{eq:open-loop-singular} satisfied, as well as
feasibility and optimality checks. Seller $B$ has the corresponding
equations with the roles of the sellers exchanged. Together these give
candidate equilibrium paths, subject to appropriate conditions at infinity;
they are not a sufficient equilibrium characterization.

\paragraph{Global verification for a moving rival.}
For a proposed rival path $u_B^*(t)$, suppose a bounded, continuously
differentiable function $U_A(t,x,y)$ satisfies, for almost every $t$,
\begin{equation}\label{eq:open-loop-verification}
\begin{aligned}
\delta U_A={}&\partial_tU_A+F_A
 +\max_{u\in[0,1]}\{\lambda D_A(p_1+u\Delta p-x)\partial_xU_A\}\\
&+\lambda D_B(p_1+u_B^*(t)\Delta p-y)\partial_yU_A.
\end{aligned}
\end{equation}
If the proposed $u_A^*(t)$ attains the maximum along the proposed state
path, it is a global best response from that initial state. Verification
for both sellers establishes open-loop Nash equilibrium;
Appendix~\ref{app:verification} gives the proof.
This is a sufficient smooth verification test; no existence or smoothness
of such functions is asserted. The maximizing rule away from the proposed
path is only a device for evaluating deviations, not a commitment by the
rival to a feedback strategy.

\section{Open questions}

\begin{enumerate}
\item Extend Proposition~\ref{prop:integer-rival} beyond integer Beta shapes,
the real-shape square $[1,2]^2$, and their neighborhoods to arbitrary real
shapes in the common-$n$ model,
or prove the weaker condition sufficient for the one-dimensional value
comparison. Integer parameter points are not dense, so continuity alone
does not complete the argument.
\item Determine when the stationary test admits an equilibrium, whether
it is unique, and whether asymmetric or boundary equilibria coexist.
\item Establish existence and characterize moving open-loop paths,
including singular arcs, switching patterns, and convergence. A rival
that always uses product~1 provides an intermediate best-response problem.
\item As a separate extension, determine how feedback responses alter
investment incentives and equilibrium existence.
\end{enumerate}

\clearpage
\appendix

\section{Model and payoff derivations}\label{app:model}

\subsection{State space and demand}

We assume $\rho\in(1/2,1)$ throughout. Then $f_0([0,1])\cup f_1([0,1])
=[0,\rho]\cup[1-\rho,1]=[0,1]$, so by induction, for every $k$, the images
of $[0,1]$ under the $2^k$ compositions of $k$ maps cover $[0,1]$, and each
image has length $\rho^k$. Every share in $[0,1]$ is therefore a limit of
shares reached along outcome histories, and the long-run state space is the
square of success shares $(x,y)\in[0,1]^2$, with $x$ for seller $A$ and $y$
for seller $B$.

\emph{Monotonicity.} For $x>x'$, the ratio of the $\Beta(1+nx,1+n(1-x))$
and $\Beta(1+nx',1+n(1-x'))$ densities is proportional to
$(\theta/(1-\theta))^{n(x-x')}$, which is increasing in $\theta$. The
posterior at $x$ therefore dominates the posterior at $x'$ in the
likelihood-ratio order, hence in the first-order stochastic order.
Conditioning on $\widetilde\theta_B$ shows that $D_A$ is nondecreasing in
$x$; the same argument shows that it is nonincreasing in $y$.

\subsection{Fluid embedding and admissibility}\label{app:fluid}

We embed the model in a family indexed by a period length $dt>0$. Each
period lasts $dt$. A sale earns $(R-c_i)\,dt$, the discount factor per
period is $\gamma^{dt}$, and an observation moves the chosen seller's share
from $x$ to $\rho^{dt}x+1-\rho^{dt}$ after a success and to $\rho^{dt}x$
after a failure. The demand \eqref{eq:boundary-demand} keeps the original
$n$. Set
\begin{equation*}
\lambda=-\log\rho,\qquad\delta=-\log\gamma,
\end{equation*}
so that $\rho^{dt}=1-\lambda\,dt+o(dt)$ and $\gamma^{dt}=1-\delta\,dt+o(dt)$.
We replace pure actions with relaxed controls $u_A,u_B\in[0,1]$, the
fraction of product~2.

If seller $A$ is chosen and uses $u_A$, its success probability is
$p_1+u_A\Delta p$, and its share changes in expectation by
\begin{equation*}
(p_1+u_A\Delta p)\lambda(1-x)\,dt-(1-p_1-u_A\Delta p)\lambda x\,dt+o(dt)
=\lambda(p_1+u_A\Delta p-x)\,dt+o(dt),
\end{equation*}
while seller $B$'s share does not move. The symmetric statement holds when
$B$ is chosen. Replacing the random motion by its expectation gives the
dynamics in \eqref{eq:duopoly-fluid}.

For any measurable control pair, the vector field is measurable in time,
uniformly Lipschitz in the state, and bounded on the square. At $x=0$ its
first component is nonnegative, and at $x=1$ it is nonpositive; the same
holds for $y$. Standard ODE existence and uniqueness and these inward
boundary signs give a unique absolutely continuous trajectory in the square.

\subsection{Payoff transformation}\label{app:transform}

\begin{proof}[Derivation of \eqref{eq:duopoly-transform}]
Since $R>c_2$, we have $q>p_1+\Delta p=p_2$. Using
$\dot x=\lambda D_A(p_1+u_A\Delta p-x)$,
\begin{align*}
D_A\bigl(R-c_1-u_A\Delta c\bigr)+\frac{\Delta c}{\lambda\Delta p}\dot x
&=D_A\Bigl[R-c_1-u_A\Delta c
+\frac{\Delta c}{\Delta p}\bigl(p_1+u_A\Delta p-x\bigr)\Bigr]\\
&=\frac{\Delta c}{\Delta p}D_A(x,y)(q-x).
\end{align*}
The product choice cancels: choosing product~2 costs $D_A\Delta c$ in flow
profit and raises the drift by $\lambda D_A\Delta p$, and the identity
credits each unit of drift at the constant rate $\Delta c/(\lambda\Delta p)$.

Along any trajectory with $x(0)=x_0$, integration by parts gives
\begin{equation*}
\int_0^T e^{-\delta t}\dot x\,dt
=e^{-\delta T}x(T)-x_0+\delta\int_0^T e^{-\delta t}x\,dt,
\end{equation*}
and the boundary term vanishes as $T\to\infty$ because $x(T)\in[0,1]$.

Substituting this identity into the discounted profit gives
\eqref{eq:duopoly-transform}.
\end{proof}

\section{Demand properties}\label{app:demand}

Unless explicitly stated otherwise, fix $n>0$ and $y\in[0,1]$, and write
\begin{equation*}
N=n+2,\qquad a=1+nx,\qquad \alpha=1+ny,\qquad \beta=1+n(1-y).
\end{equation*}
Seller $A$'s posterior is $\Beta(a,N-a)$ and the rival's is
$\Beta(\alpha,\beta)$, with density $f_B$ and distribution function $F_B$.
Since $\partial_x=n\partial_a$, log-concavity in $x$ and $a$ is equivalent.

\subsection{Integer and real rival shapes}\label{app:integer-rival}

\begin{proof}[Proof of Proposition~\ref{prop:integer-rival}: integer shapes]
Set $M=\alpha+\beta-1$ and write $(v)_j=v(v+1)\cdots(v+j-1)$,
with $(v)_0=1$. The integer-shape Beta distribution has CDF
\begin{equation*}
F_{\alpha,\beta}(p)=\sum_{j=\alpha}^{M}\binom{M}{j}
p^j(1-p)^{M-j}.
\end{equation*}
Indeed, the sum is zero at $p=0$, and differentiating it telescopes to the
$\Beta(\alpha,\beta)$ density. Taking the expectation under
$\Beta(a,N-a)$ and evaluating its mixed moments yields
\begin{equation}\label{eq:integer-demand-polynomial}
P(a):=\mathcal D(a)=\frac{1}{(N)_M}
\sum_{j=\alpha}^{M}\binom{M}{j}(a)_j(N-a)_{M-j}.
\end{equation}
The right-hand side is a polynomial of degree at most $M$. Similarly,
\begin{equation*}
1-P(a)=\frac{1}{(N)_M}
\sum_{j=0}^{\alpha-1}\binom{M}{j}(a)_j(N-a)_{M-j}.
\end{equation*}
These identities initially hold on $(0,N)$ and hence hold as polynomial
identities on the real line. The first shows that $(a)_\alpha$ divides
$P(a)$; the second shows that $(N-a)_\beta$ divides $1-P(a)$.
Consequently,
\begin{equation*}
P(0)=P(-1)=\cdots=P(1-\alpha)=0,
\qquad
P(N)=P(N+1)=\cdots=P(N+\beta-1)=1.
\end{equation*}

Rolle's theorem gives a zero $r_j$ of $P'$ in each $(-j,1-j)$,
$j=1,\ldots,\alpha-1$, and a zero $s_j$ in each
$(N+j-1,N+j)$, $j=1,\ldots,\beta-1$. These are $M-1$ distinct real
zeros. Since $P'$ is nonzero and has degree at most $M-1$, they exhaust
its zeros, which are all simple and outside $(0,N)$. In particular, $P'$
has constant sign there; $P(N)-P(0)=1$ makes this sign positive.
Factoring $P'$ now gives
\begin{equation}\label{eq:integer-curvature}
\frac{d^2}{da^2}\log\mathcal D'(a)
=-\sum_{j=1}^{\alpha-1}\frac{1}{(a-r_j)^2}
 -\sum_{j=1}^{\beta-1}\frac{1}{(s_j-a)^2}<0
\end{equation}
whenever $M>1$. If $M=1$, then $P(a)=a/N$ and $P'=1/N$.
\end{proof}

The root locations also give the explicit bound, for $\alpha+\beta>2$,
\begin{equation}\label{eq:integer-curvature-bound}
\frac{d^2}{da^2}\log\mathcal D'(a)
< -\sum_{j=1}^{\alpha-1}\frac{1}{(a+j)^2}
  -\sum_{j=1}^{\beta-1}\frac{1}{(N-a+j)^2}.
\end{equation}
Thus the result concerns every own reputation, not just integer values of
$a$. Nor does it require $N=\alpha+\beta$. For example, against
$\Beta(2,2)$ with $N=4$,
\begin{equation*}
\mathcal D(a)=\frac{a(a+1)(7-a)}{60},\qquad
\mathcal D'(a)=\frac{-3a^2+12a+7}{60},
\end{equation*}
whose roots are $2\pm\sqrt{57}/3$, in $(-1,0)$ and $(4,5)$.

\paragraph{A real-parameter extension.}
The second part of Proposition~\ref{prop:integer-rival} follows from a
shape-preservation argument that applies to any concave rival density.

\begin{lemma}[Concave rival density]\label{lem:concave-rival-density}
Let $N>0$ and let $F$ be a distribution function on $[0,1]$ with positive
density $f\in C^2(0,1)$. If $f''\leq-c$ for a constant $c\geq0$, then
\begin{equation*}
H(a)=\E[F(X_a)],\qquad X_a\sim\Beta(a,N-a),\qquad 0<a<N,
\end{equation*}
satisfies
\begin{equation}\label{eq:concave-density-bound}
H'>0,\qquad H'''\leq-\frac{c}{N(N+1)(N+2)}.
\end{equation}
Thus $H'$ is concave and log-concave, with strictly negative log curvature
if $c>0$.
\end{lemma}

\begin{proof}
Fix $a\in(0,N)$, write $b=N-a$, and take $0<h<b/3$. Couple the four
posterior laws by
\begin{equation*}
(U,V_1,V_2,V_3,W)\sim\operatorname{Dirichlet}(a,h,h,h,b-3h).
\end{equation*}
For every subset $A\subseteq\{1,2,3\}$, the variable
$U+\sum_{i\in A}V_i$ has distribution $\Beta(a+|A|h,b-|A|h)$.
Writing $\Delta_vF(u)=F(u+v)-F(u)$, it follows that
\begin{align*}
\Delta_h^3 H(a)
&=\E[\Delta_{V_1}\Delta_{V_2}\Delta_{V_3}F(U)]\\*
&=\E\left[\int_0^{V_1}\int_0^{V_2}\int_0^{V_3}
f''(U+t_1+t_2+t_3)\,dt_3\,dt_2\,dt_1\right]\\*
&\leq-c\,\E[V_1V_2V_3]
=-\frac{ch^3}{N(N+1)(N+2)}.
\end{align*}
All integration intervals lie strictly inside $(0,1)$ almost surely.
Also, $H$ is smooth in $a$: since $F$ is bounded, derivatives of its
Beta integral are dominated on compact parameter intervals by integrable
powers times logarithmic factors. Dividing by $h^3$ and letting
$h\downarrow0$ gives the bound on $H'''$.
Finally, $H'(a)>0$ follows by integrating the positive derivative
\eqref{eq:dirichlet-derivative} against $f$ (that identity only needs
$a,b>0$). Therefore
\begin{equation*}
(\log H')''=\frac{H'''}{H'}-\left(\frac{H''}{H'}\right)^2,
\end{equation*}
which proves both log-concavity claims.
\end{proof}

\begin{proof}[Proof of Proposition~\ref{prop:integer-rival}: the real-shape range]
Set $r=\alpha-1$ and $s=\beta-1$. For $r,s\in[0,1]$, the rival density
satisfies
\begin{equation*}
\frac{f''(p)}{f(p)}
=\frac{r(r-1)}{p^2}+\frac{s(s-1)}{(1-p)^2}
 -\frac{2rs}{p(1-p)}\leq0.
\end{equation*}
Except for the three affine densities $(\alpha,\beta)=(1,1),(2,1),(1,2)$,
we have the uniform bound $f''\leq-c<0$, where one may take
\begin{equation*}
c=\frac{1}{\Bfun(\alpha,\beta)}
\begin{cases}
2rs,&r,s>0,\\
r(1-r),&s=0,\quad0<r<1,\\
s(1-s),&r=0,\quad0<s<1.
\end{cases}
\end{equation*}
For the first case use $p^{r-1}(1-p)^{s-1}\geq1$ in the cross term;
the boundary cases use $p^{r-2}\geq1$ or $(1-p)^{s-2}\geq1$.
Lemma~\ref{lem:concave-rival-density} therefore gives
\begin{equation*}
(\log\mathcal D')''\leq
-\frac{c}{N(N+1)(N+2)\mathcal D'}
-\left(\frac{\mathcal D''}{\mathcal D'}\right)^2<0.
\end{equation*}
For the two nonconstant affine densities, direct Beta moments give
\begin{equation*}
\mathcal D'_{2,1}(a)=\frac{2a+1}{N(N+1)},\qquad
\mathcal D'_{1,2}(a)=\frac{2(N-a)+1}{N(N+1)},
\end{equation*}
whose logarithms have strictly negative second derivatives. For $(1,1)$,
$\mathcal D(a)=a/N$ as before.
\end{proof}

\begin{corollary}[Neighborhoods of proved parameters]
\label{cor:integer-neighborhood}
Let $(n_0,y_0)\in(0,\infty)\times[0,1]$ have rival shapes that are either
both integers or both in $[1,2]$. There is a relative neighborhood of
$(n_0,y_0)$ in $(0,\infty)\times[0,1]$ on which
$x\mapsto\partial_xD_A(x,y)$ is strictly log-concave on $[0,1]$.
\end{corollary}

\begin{proof}
At $(n_0,y_0)$, the rival shapes sum to $n_0+2>2$.
Proposition~\ref{prop:integer-rival} gives strictly negative log curvature
for every $x\in[0,1]$, with $\partial_xD_A>0$. These quantities are
continuous in $(n,y,x)$: on a compact parameter neighborhood the Beta
integrals and their first three $x$-derivatives admit integrable bounds,
including the logarithmic factors introduced by differentiation.
Compactness of $[0,1]$ therefore gives a uniform negative curvature
margin and a uniform positive lower bound on marginal demand. Both
persist in a sufficiently small parameter neighborhood.
\end{proof}

The real-shape argument proves a continuous parameter range directly;
it does not interpolate between integer points. It still does not settle
the general conjecture: for shapes at least one, a Beta density ceases to
be concave when either shape exceeds two. Integer shape pairs are not
dense, and the continuity argument provides no lower bound on neighborhood
sizes that would cover all remaining parameters.

\subsection{Log-concave marginal demand against a known benchmark}

\begin{lemma}[Single-seller marginal demand]
\label{lem:single-seller-logconcave}
Let $N>2$, $p\in(0,1)$, and $S(a;p)=\Pb(\Beta(a,N-a)>p)$ for $0<a<N$. Then
\begin{equation}\label{eq:dirichlet-derivative}
\partial_aS(a;p)=\frac{1}{\Bfun(a,N-a)}
\int_0^p\int_0^{1-p}\frac{u^{a-1}w^{N-a-1}}{1-u-w}\,dw\,du>0,
\end{equation}
and $a\mapsto\partial_aS(a;p)$ is strictly log-concave on $[1,N-1]$. In
particular, with $N=n+2$, $S_{p_0}'(x)=n\,\partial_aS(1+nx;p_0)$ is positive
and strictly log-concave on $[0,1]$.
\end{lemma}

\begin{proof}
Write $b=N-a$. For $0<h<b$, let $(U,Y,W)$ be Dirichlet with parameters
$(a,h,b-h)$. Then $U\sim\Beta(a,b)$ and $U+Y=1-W\sim\Beta(a+h,b-h)$, and
$U+Y\geq U$, so
\begin{equation*}
S(a+h;p)-S(a;p)=\Pb(U\leq p<U+Y)=\Pb(U<p,\ W<1-p).
\end{equation*}
Integrating the Dirichlet density over this rectangle, dividing by $h$, and
letting $h\downarrow0$ (using $h\Gamma(h)\to1$) gives
\eqref{eq:dirichlet-derivative}. Dominated convergence applies because the
singularity of $(1-u-w)^{-1}$ at the corner $(p,1-p)$ is integrable in two
dimensions.

Expanding $(1-u-w)^{-1}=\sum_{m\geq0}(u+w)^m$ and integrating term by term
(Tonelli),
\begin{equation*}
\partial_aS=\sum_{j,k\geq0}t_{jk}(a),\qquad
t_{jk}(a)=\frac{p^a(1-p)^b}{\Bfun(a,b)}
\binom{j+k}{j}\frac{p^j(1-p)^k}{(a+j)(b+k)}.
\end{equation*}
Differentiating in $a$ with $a+b=N$ fixed,
\begin{align*}
(\log t_{jk})'&=\log\tfrac{p}{1-p}-\psi(a)+\psi(b)
-\tfrac1{a+j}+\tfrac1{b+k},\\
(\log t_{jk})''&=-\psi_1(a)-\psi_1(b)+\tfrac1{(a+j)^2}+\tfrac1{(b+k)^2},
\end{align*}
where $\psi$ is the digamma function. For a positive sum, with expectations
and variances taken under the weights $t_{jk}/\partial_aS$,
\begin{equation*}
(\log\partial_aS)''
=\E\bigl[(\log t_{jk})''\bigr]+\Var\bigl((\log t_{jk})'\bigr).
\end{equation*}
Termwise differentiation is justified on compact subintervals of $(0,N)$,
where the summands and their first two derivatives are bounded by a constant
times the summands at a fixed interior parameter. The random part of
$(\log t_{jk})'$ is $-1/(a+j)+1/(b+k)\in[-1/a,1/b]$, so its variance is at
most $(1/a+1/b)^2/4$. Therefore
\begin{equation*}
(\log\partial_aS)''\leq-\psi_1(a)-\psi_1(b)+\frac1{a^2}+\frac1{b^2}
+\frac14\Bigl(\frac1a+\frac1b\Bigr)^2.
\end{equation*}
For $s>0$, $\psi_1(s)=\sum_{m\geq0}(s+m)^{-2}>1/s+1/(2s^2)$, by summing the
strict trapezoidal inequality for the convex function $t\mapsto(s+t)^{-2}$
over $[m,m+1]$. With $A=1/a$ and $B=1/b$, both in $(0,1]$ when
$a,b\geq1$,
\begin{equation*}
(\log\partial_aS)''<-(A+B)+\tfrac34(A^2+B^2)+\tfrac12AB
\leq-\tfrac14(A^2+B^2)+\tfrac12AB=-\tfrac14(A-B)^2\leq0,
\end{equation*}
where the second inequality uses $A+B\geq A^2+B^2$. The final statement
follows from $S_{p_0}(x)=S(1+nx;p_0)$, since $a=1+nx$ and $b=1+n(1-x)$ are
at least one on $[0,1]$.
\end{proof}

\subsection{An exact representation of marginal demand}

\begin{lemma}[Marginal demand as a Beta smoothing]
\label{lem:marginal-demand-representation}
Define, for $\pi\in(0,1)$,
\begin{equation*}
k_y(\pi)=\int_0^1(1-h)^{N-1}\int_0^1
f_B\bigl((1-h)\pi+ht\bigr)\,dt\,dh.
\end{equation*}
Then, with $\Pi_a\sim\Beta(a,N-a)$,
\begin{equation}\label{eq:marginal-demand-representation}
\partial_aD_A=\E\bigl[k_y(\Pi_a)\bigr]>0.
\end{equation}
Moreover, writing $T=\log\bigl(\Pi/(1-\Pi)\bigr)$, $P_a$ for the
$\Beta(a,N-a)$ law, and $Q_a$ for the law with density proportional to
$k_y$ with respect to $P_a$,
\begin{equation}\label{eq:logconcavity-variance}
\frac{d^2}{da^2}\log\partial_aD_A
=\Var_{Q_a}(T)-\Var_{P_a}(T)
=\Var_{Q_a}(T)-\bigl[\psi_1(a)+\psi_1(N-a)\bigr],
\end{equation}
where $\psi_1$ is the trigamma function.
\end{lemma}

\begin{proof}
Let $S(a;p)=\Pb(\Beta(a,N-a)>p)$. By \eqref{eq:dirichlet-derivative} in
Lemma~\ref{lem:single-seller-logconcave}, for $p\in(0,1)$,
\begin{equation*}
\partial_aS(a;p)=\frac{1}{\Bfun(a,N-a)}
\int_0^p\int_0^{1-p}
\frac{u^{a-1}w^{N-a-1}}{1-u-w}\,dw\,du.
\end{equation*}
Since $D_A=\int_0^1f_B(p)S(a;p)\,dp$, differentiating under the integral and
applying Fubini's theorem give
\begin{equation*}
\partial_aD_A=\frac{1}{\Bfun(a,N-a)}
\iint_{u,w>0,\,u+w<1}u^{a-1}w^{N-a-1}\,
\frac{F_B(1-w)-F_B(u)}{1-u-w}\,du\,dw.
\end{equation*}
Change variables to $u=(1-h)\pi$ and $w=(1-h)(1-\pi)$, with
$\pi,h\in(0,1)$ and Jacobian $1-h$. Then
$u^{a-1}w^{N-a-1}\,du\,dw=(1-h)^{N-1}\pi^{a-1}(1-\pi)^{N-a-1}\,d\pi\,dh$,
and since $1-w=(1-h)\pi+h$ and $1-u-w=h$,
\begin{equation*}
\frac{F_B(1-w)-F_B(u)}{1-u-w}
=\frac1h\int_{(1-h)\pi}^{(1-h)\pi+h}f_B
=\int_0^1f_B\bigl((1-h)\pi+ht\bigr)\,dt.
\end{equation*}
This proves \eqref{eq:marginal-demand-representation}. Positivity follows
because $f_B>0$ on $(0,1)$.

For the second claim, $\pi^{a-1}(1-\pi)^{N-a-1}=e^{aT}\pi^{-1}(1-\pi)^{N-1}$.
Hence
\begin{equation*}
\log\partial_aD_A=\log\int_0^1 e^{aT}k_y(\pi)\pi^{-1}(1-\pi)^{N-1}\,d\pi
-\log\Bfun(a,N-a).
\end{equation*}
The second derivative of the first term is the variance of $T$ under the
tilted law $Q_a$. The second derivative of $\log\Bfun(a,N-a)$ is
$\psi_1(a)+\psi_1(N-a)$, which is also $\Var_{P_a}(T)$, because
$T=\log G_a-\log G_{N-a}$ for independent gamma variables
$G_a\sim\mathrm{Gamma}(a)$ and $G_{N-a}\sim\mathrm{Gamma}(N-a)$.
\end{proof}

\paragraph{Reading of the lemma.}
The function $k_y$ does not depend on $a$: it is a smoothing of the rival's
density, evaluated along short random segments around $\pi$. Marginal
demand is the expectation of $k_y$ under one's own posterior. By
\eqref{eq:logconcavity-variance}, log-concavity of marginal demand is
equivalent to the statement that \emph{reweighting one's own posterior by
$k_y$ does not increase the variance of its logit}. Two boundary cases
confirm the formula.
\begin{itemize}
\item \emph{Uninformative rival} ($f_B\equiv1$): $k_y\equiv1/N$, so
$Q_a=P_a$ and the second derivative is zero. Indeed $D_A=a/N$ is linear.
\item \emph{Point-mass rival} ($f_B=\delta_{p_0}$): $k_y$ localizes $\Pi$
near $p_0$, and \eqref{eq:marginal-demand-representation} reduces to the
single-seller marginal demand, which is strictly log-concave by
Lemma~\ref{lem:single-seller-logconcave}.
\end{itemize}
Log-concave reweighting does not reduce variance in general, and a direct
Pr\'ekopa argument fails: $k_y(0^+)>0$, so the joint density of
$(a,\pi)$ is not log-concave near $\pi=0$. A proof must therefore use the
specific structure of $k_y$.

\paragraph{An obstruction to a generic preservation argument.}
The smoothing kernel itself need not be log-concave, even at parameters
covered by Proposition~\ref{prop:integer-rival}. For $N=4$ and a
$\Beta(1,3)$ rival, evaluating its defining integral gives
\begin{equation*}
k_y(\pi)=\frac{(1-\pi)^2}{2}+\frac{1-\pi}{10}+\frac1{60},
\qquad
\left.\frac{d^2}{d\pi^2}\log k_y(\pi)\right|_{\pi=1}=24>0.
\end{equation*}
Thus log-concavity of the rival density does not imply log-concavity of
$k_y$. The Beta-operator preservation result of \cite{badia2023}, which
assumes a monotone log-concave input, cannot be applied directly to this
kernel. The polynomial proof above avoids that requirement.

\subsection{Exploratory numerical checks}\label{app:demand-numerics}

For a fixed $(N,\alpha,\beta)$, the derivatives of demand can be evaluated
without finite differences. Let $P\sim\Beta(a,b)$, $b=N-a$, and define
\begin{equation*}
Z=\log\frac{P}{1-P}-\psi(a)+\psi(b),\qquad
v=\psi_1(a)+\psi_1(b),\qquad
w=\psi_2(a)-\psi_2(b).
\end{equation*}
Differentiating the density under the integral gives
\begin{equation*}
\begin{aligned}
\mathcal D'(a)&=\E[F_B(P)Z],\\
\mathcal D''(a)&=\E[F_B(P)(Z^2-v)],\\
\mathcal D'''(a)&=\E[F_B(P)(Z^3-3vZ-w)].
\end{aligned}
\end{equation*}
The diagnostic is $\mathcal D'''/\mathcal D'-(\mathcal D''/\mathcal D')^2$,
the curvature of $\log\mathcal D'$. In the common-$n$ model, its sign is
the same as the curvature with respect to $x$, which is $n^2$ times larger.

The accompanying scripts evaluate these integrals in logit coordinates,
rescale their integrands to handle very small demand, and use symmetry to
avoid computing derivatives from a demand close to one. An exploratory
scan used
\begin{equation*}
\begin{aligned}
n&\in\{0.01,0.1,0.5,1,2,3,5,10,20,50,100,300\},\\
x,y&\in\{0,0.01,0.1,0.25,0.5,0.75,0.9,0.99,1\}.
\end{aligned}
\end{equation*}
All $972$ computed curvatures were negative; the largest was approximately
$-0.00165906$, at $n=300$ and $x=y=0$ (equivalently $x=y=1$).
There are $567$ points with the economic restriction $n>2$; the other
points test the broader mathematical conjecture.

Independent checks compare the score integrals with the exact polynomial
structure and with high-precision integration in the original probability
coordinate. Additional checks for real shapes in $[1,2]^2$ compare the
third-derivative score formula with centered differences of the original
probability integral, including unequal posterior totals and the three
affine-density exceptions. These are floating-point diagnostics, not interval certificates
or a proof between grid points. The scripts and their outputs are in
\texttt{Progress notes/demand\_shape/}.

\subsection{The demand jump against a known benchmark}

\begin{lemma}[Jump identity]\label{lem:jump-identity}
For $a,b>0$ and $u\in(0,1)$,
\begin{equation*}
\Pb\bigl(\Beta(a+1,b)>u\bigr)-\Pb\bigl(\Beta(a,b+1)>u\bigr)
=\frac{\Gamma(a+b+1)}{\Gamma(a+1)\Gamma(b+1)}u^a(1-u)^b
=\frac{f_{a+1,b+1}(u)}{a+b+1},
\end{equation*}
where $f_{a+1,b+1}$ is the $\Beta(a+1,b+1)$ density.
\end{lemma}

\begin{proof}
Let $I_u(a,b)=\Pb(\Beta(a,b)\leq u)$. Integration by parts in the defining
integrals, together with $\Bfun(a+1,b)=a\Bfun(a,b)/(a+b)$ and
$\Bfun(a,b+1)=b\Bfun(a,b)/(a+b)$, gives the contiguous relations
\begin{equation*}
I_u(a+1,b)=I_u(a,b)-\frac{u^a(1-u)^b}{a\Bfun(a,b)},
\qquad
I_u(a,b+1)=I_u(a,b)+\frac{u^a(1-u)^b}{b\Bfun(a,b)}.
\end{equation*}
The left-hand side of the lemma is $I_u(a,b+1)-I_u(a+1,b)
=u^a(1-u)^b(a+b)/(ab\Bfun(a,b))$. Using
$\Bfun(a,b)=\Gamma(a)\Gamma(b)/\Gamma(a+b)$ and $s\Gamma(s)=\Gamma(s+1)$
gives both expressions.
\end{proof}

\subsection{The finite-observation analogue and concentrated rivals}

For shape parameters $(a,b)$ of seller $A$'s posterior, write
$D_A(a,b)=\Pb(\Beta(a,b)>\widetilde\theta_B)$.

\begin{proposition}[The per-observation analogue is exactly log-concave]
\label{prop:discrete-jump-logconcave}
Fix a pre-update total $M>0$ and let seller $A$'s shape parameters be
$(a,M-a)$, $a\in(0,M)$. The demand jump from a success instead of a failure,
\begin{equation*}
J(a)=D_A(a+1,M-a)-D_A(a,M-a+1),
\end{equation*}
is log-concave in $a$, strictly unless $\alpha=\beta=1$.
\end{proposition}

\begin{proof}
Condition on the rival's draw $\widetilde\theta_B=u$ and apply
Lemma~\ref{lem:jump-identity} with $b=M-a$. Integrating against $f_B$,
\begin{equation*}
J(a)=\frac{1}{M+1}\int_0^1f_{a+1,M-a+1}f_B
=\frac{\Bfun(a+\alpha,M-a+\beta)}
{(M+1)\,\Bfun(a+1,M-a+1)\,\Bfun(\alpha,\beta)},
\end{equation*}
where $f_{a+1,M-a+1}$ is the $\Beta(a+1,M-a+1)$ density. The terms
$\Gamma(M+\alpha+\beta)$ and $\Gamma(M+2)$ do not depend on $a$, so
\begin{equation*}
\frac{d^2}{da^2}\log J
=\bigl[\psi_1(a+\alpha)-\psi_1(a+1)\bigr]
+\bigl[\psi_1(M-a+\beta)-\psi_1(M-a+1)\bigr]\leq0,
\end{equation*}
since $\psi_1$ is strictly decreasing and $\alpha,\beta\geq1$. Equality
requires $\alpha=\beta=1$.
\end{proof}

The jump $J(a)$ is the integral of the marginal demand
$\partial_tD_A(t,M+1-t)$ over $t\in[a,a+1]$. It is the quantity that enters
the Bellman equation of the discrete model, and its log-concavity is exact
for every rival. Log-concavity of $\partial_aD_A$ implies log-concavity of
these unit-window averages, but not conversely. The difficulty is therefore
specific to the infinitesimal step of the fluid limit.

\begin{proposition}[Concentrated rival]\label{prop:concentrated-rival}
Fix $N$ and $p_0\in(0,1)$, and let the rival's posterior be any law $\mu$ on
$[0,1]$. There is a weak neighborhood $\mathcal U$ of $\delta_{p_0}$ such
that, for every $\mu\in\mathcal U$, $x\mapsto\partial_xD_A(x;\mu)$ is
strictly log-concave on $[0,1]$. In particular, the conclusion of
Proposition~\ref{prop:frozen-reduction} holds when the rival's reputation is
known much more precisely than one's own, for example under asymmetric
forgetting rates with $n_B\gg n_A$.
\end{proposition}

\begin{proof}[Proof sketch]
Write $D(x;\mu)=\int S_p(x)\,\mu(dp)$. Since $a,N-a\geq1$, $S_p(x)$ and its
first three $x$-derivatives are jointly continuous in $(x,p)\in[0,1]^2$:
the $a$-derivatives of the incomplete beta integral only bring in factors
$\log t$, which are integrable against $t^{a-1}$. By uniform continuity on a
compact set, $\partial_x^jD(\cdot;\mu)\to\partial_x^jS_{p_0}$ uniformly on
$[0,1]$ as $\mu\to\delta_{p_0}$ weakly, for $j\leq3$. By
Lemma~\ref{lem:single-seller-logconcave}, $S_{p_0}'>0$ and
$(\log S_{p_0}')''<0$ on $[0,1]$; by continuity and compactness,
$S_{p_0}'\geq c>0$ and $(\log S_{p_0}')''\leq-\eta<0$. Since
$(\log D')''=(D'D'''-(D'')^2)/(D')^2$ depends continuously on
$(D',D'',D''')$ where $D'\geq c/2$, it stays below $-\eta/2$ for $\mu$ close
enough to $\delta_{p_0}$.
\end{proof}

\subsection{Large-\texorpdfstring{$n$}{n} heuristics}

\begin{remark}[Large-$n$ heuristic]\label{rem:large-n}
For large $n$, each posterior has variance approximately $x(1-x)/n$, so
$D(x)\approx\Phi(z)$ with $z=(x-y)\sqrt{n/M(x)}$ and
$M(x)=x(1-x)+y(1-y)$. Then $\log D'\approx-\tfrac n2g(x)+O(\log n)$, where
$g=(x-y)^2/M$, and the remaining term $\log\partial_xz$ has $O(1)$ second
derivative. A direct calculation gives
\begin{equation*}
\tfrac12M^3g''=\bigl[M-(x-y)(1-2x)\bigr]^2+(x-y)^2M>0,
\end{equation*}
which is bounded below on $[0,1]$ for each $y\in(0,1)$, because
$M\geq y(1-y)$. The leading term is therefore strictly concave with
curvature of order $n$, and dominates for large $n$. This suggests that
log-concavity holds for $n$ large, with a threshold that deteriorates as $y$
approaches $0$ or $1$, where the approximation also breaks down. With
$y(1-y)$ replaced by $0$, the same identity is the leading-order version of
Lemma~\ref{lem:single-seller-logconcave}. This is a heuristic, not a proof.
\end{remark}

\paragraph{Heuristic scaling.}
For large $n$, the normal approximation of Remark~\ref{rem:large-n} gives
\begin{equation*}
\partial_xD_A(x,x)\approx\sqrt{\frac{n}{4\pi x(1-x)}}.
\end{equation*}
Compare this with the single-seller benchmark:
\begin{equation*}
S_{p_0}'(x)\approx\sqrt{\frac{n}{x(1-x)}}\;
\phi\!\left((x-p_0)\sqrt{\frac{n}{x(1-x)}}\right),
\end{equation*}
which is exponentially small in $n$ unless $x$ is within $O(n^{-1/2})$ of
$p_0$. In the duopoly, the benchmark is the rival's reputation, which at a
symmetric state sits where one's own does. Competition keeps the customer
close to indifference, and the slope of demand is of order $\sqrt n$
everywhere on the diagonal.

As $\rho\uparrow1$, $\lambda\sim1-\rho$ and $n=(1-\rho)^{-1}$, so the
left-hand side of \eqref{eq:stationary-peak-test} grows like $(1-\rho)^{-1/2}$
and the right-hand side like $\delta/(1-\rho)$. The condition eventually
fails, so interior steady states disappear under slow forgetting. However,
the duopoly should sustain product~2 for a much wider range of forgetting
rates than the single-seller benchmark. This is heuristic and is not yet
supported by a proof or by computations.

\section{Frozen-rival control and the single-interval result}\label{app:control}

\subsection{One-dimensional control with a valley-peak payoff}

Let $D$ be smooth and positive on $[0,1]$, let
$b_i(x)=\lambda D(x)(p_i-x)$ for $i\in\{1,2\}$, and consider
\begin{equation}\label{eq:1d-problem}
W(x)=\sup\int_0^\infty e^{-\delta t}F(x(t))\,dt
\quad\text{s.t.}\quad
\dot x=(1-u)b_1(x)+ub_2(x),\ u(t)\in[0,1],\ x(0)=x.
\end{equation}
Note that $b_1<b_2$ everywhere; both are positive below $p_1$, $b_1<0<b_2$
on $(p_1,p_2)$, and both are negative above $p_2$. Every state in
$[p_1,p_2]$ can be held by the mixture $u=(x-p_1)/\Delta p$. The value
satisfies the HJB equation $\delta W=F+\max_i b_iW'$, so product~2 is
optimal where $W'>0$ and product~1 where $W'<0$.

\begin{lemma}[At most one interval of product~2]\label{lem:control}
Let $F\in C^2[0,1]$. Suppose that either $F$ is nonincreasing, or there are
$0\leq a<m\leq1$ with $F'<0$ on $[0,a)\cup(m,1]$, $F'>0$ on $(a,m)$,
$F''(a)>0$ if $a>0$, and $F''(m)<0$ if $m<1$. Then \eqref{eq:1d-problem}
admits an optimal stationary policy that uses product~2 on at most one
interval, with a holding mixture at $m$ if $m\in(p_1,p_2)$ is an attractive
threshold.
\end{lemma}

\begin{proof}
\emph{Step 0: monotone cases.} Let $x_i(t;x)$ be the trajectory under
product $i$. By the comparison theorem for scalar ODEs, every admissible
trajectory from $x$ satisfies $x_1(t;x)\leq x(t)\leq x_2(t;x)$. If $F$ is
nonincreasing, product~1 therefore maximizes $F(x(t))$ at every $t$ and is
optimal everywhere. If $F$ is nondecreasing ($a=0$, $m=1$), product~2 is
optimal everywhere. Assume from now on that $0\leq a<m\leq1$ and the shape
is not monotone.

\emph{Step 1: switches outside $[p_1,p_2]$.} Let $g=\delta W-F
=\max_ib_iW'$. Below $p_1$ both drifts are positive, so $g$ has the sign of
$W'$, and product~2 is optimal where $g>0$. Above $p_2$ both are negative,
so $g$ has the sign of $-W'$, and product~2 is optimal where $g<0$. On
either region the HJB equation can be solved for $W'$ as
$W'=\Phi(x,W)$ with $\Phi$ locally Lipschitz (the two branches
$(\delta W-F)/b_i$ coincide where $\delta W=F$), so $W$ is $C^1$ there. On a
branch with product $i$, differentiating $\delta W=F+b_iW'$ gives
\begin{equation*}
g'=\frac{\delta g}{b_i}-F'.
\end{equation*}
At a switch point $g=0$, so $g'=-F'$. When $F'\neq0$ this fixes the
direction of the switch, in increasing $x$:
\begin{equation*}\tag{$*$}
\begin{array}{c|cc}
 &F'<0&F'>0\\ \hline
x<p_1&1\to2&2\to1\\
x>p_2&2\to1&1\to2
\end{array}
\end{equation*}
If $g=0$ at a zero of $F'$, then $g'=0$ and $g''=-F''\neq0$, so $g$ touches
zero without changing sign and no switch occurs.

\emph{Step 2: the comparison inside $[p_1,h]$.} Suppose $m>p_1$ and put
$h=\min\{m,p_2\}$. Let $L(x)$ be the value of using product~1 forever, which
drives the state down to $p_1$. Let $H(x)$ be the value of using product~2
until reaching $h$ and then holding there, or, if $h=p_2$, of using
product~2 forever. On $(p_1,h)$,
\begin{equation*}
\delta L=F+b_1L',\qquad \delta H=F+b_2H',
\qquad L(p_1)=\frac{F(p_1)}{\delta},\qquad H(h)=\frac{F(h)}{\delta}.
\end{equation*}
\begin{enumerate}
\item[(i)] $\delta\max\{L,H\}>F$ on $(p_1,h)$. If $x<a$, $F$ increases along
the descending product-1 path, so $L(x)>F(x)/\delta$. If $x\geq a$, $F$
increases along the ascending product-2 path to $h$, so $H(x)>F(x)/\delta$.
\item[(ii)] Where $L>H$, $L'=(\delta L-F)/b_1<0$, so product~1 attains
$\max_ib_iL'$. Where $H>L$, $H'=(\delta H-F)/b_2>0$, so product~2 attains
$\max_ib_iH'$.
\item[(iii)] At a crossing $L(x^*)=H(x^*)$,
$(H-L)'(x^*)=(\delta L-F)(1/b_2-1/b_1)>0$. Every crossing has the same
orientation, so there is at most one, with $L$ winning to its left and $H$
to its right.
\item[(iv)] $W=\max\{L,H\}$ on $[p_1,h]$. No trajectory starting in
$[p_1,h]$ can go below $p_1$, since $b_1(p_1)=0$. If $h=m<p_2$, an
excursion above $h$ can be replaced by holding at $h$ until the trajectory
returns, which weakly increases the payoff because $F\leq F(h)$ on $[h,1]$.
On $[p_1,h]$, $\widehat W=\max\{L,H\}$ satisfies
$\tfrac{d}{dt}\widehat W(x(t))\leq\delta\widehat W-F$ along every admissible
trajectory, for almost every $t$: by (ii) at points where $L\neq H$, and by
(i) at times spent at the crossing or an endpoint. Integrating
$\tfrac{d}{dt}[e^{-\delta t}\widehat W(x(t))]\leq-e^{-\delta t}F(x(t))$
gives $W\leq\widehat W$. Conversely $W\geq L$ and $W\geq H$, because these
values are attained by admissible policies.
\end{enumerate}
Hence, on $[p_1,h]$, any switch is $1\to2$ and occurs at most once, and
following the winning branch keeps the state in the winning region.

\emph{Step 3: assembly.}

\emph{Case $m\leq p_1$.} From $x\geq m$, every trajectory stays in $[m,1]$,
where $F$ is decreasing, so product~1 (the smallest trajectory) is optimal.
Below $m$, table $(*)$ allows only a $1\to2$ switch on $[0,a)$ and only a
$2\to1$ switch on $(a,m)$, with none at $a$. Hence product~2 is used on at
most one interval.

\emph{Case $p_1<m\leq p_2$.} Above $m$, product~1 is optimal: its trajectory
is the smallest, so it reaches $m$ first while $F$ decreases above $m$, and
it can hold at $m$ (with a mixture, or with product~2 if $m=p_2$) until any
other trajectory arrives. On $[p_1,m]$, Step~2 gives at most one $1\to2$
switch. Below $p_1$, product~2 cannot form a detached interval:
\begin{itemize}
\item If $a>p_1$, then $F'<0$ below $p_1$, and table $(*)$ allows only
$1\to2$ switches there. If $L$ wins at $p_1$, then $g(p_1)=0$; product~2
just left of $p_1$ would require $g>0$ with $g'=\delta g/b_2-F'>0$, which is
incompatible with $g$ decreasing to $0$. So product~1 is used just left of
$p_1$ and on all of $[0,p_1]$. If $H$ wins strictly at $p_1$, it wins on all
of $[p_1,m)$ by (iii), so any product-2 region below $p_1$ attaches to it.
\item If $a<p_1$, then for $x\in(a,p_1)$ the state moves up toward $p_1$
under either product and $F$ increases, so $W>F/\delta$, $g>0$, and
product~2 is used on $(a,p_1)$. Likewise $H(p_1)>F(p_1)/\delta=L(p_1)$, so
$H$ wins on $[p_1,m)$. Below $a$, only a $1\to2$ switch is possible, so the
product-2 region is a single interval ending at $m$.
\item If $a=p_1$, then moving up strictly improves on holding at the valley,
so $H$ wins at $p_1$ and the previous argument applies.
\item If $L(p_1)=H(p_1)$, then $F'(p_1)<0$, $H'(p_1)=0$, and
$L'(p_1)=F'(p_1)/(\delta+\lambda D(p_1))<0$, so $H$ wins just right of
$p_1$, while the $g'$ argument above gives product~1 just left. This is a
single $1\to2$ switch.
\end{itemize}

\emph{Case $m>p_2$.} Step~2 applies on $[p_1,p_2]$ with $h=p_2$, and the
argument below $p_1$ is unchanged.
\begin{itemize}
\item If $a\geq p_2$, every trajectory starting at $x\leq a$ stays below
$a$, where $F$ is decreasing, so product~1 is optimal there. Above $a$,
table $(*)$ allows only a $1\to2$ switch on $(a,m)$ and only a $2\to1$
switch on $(m,1]$.
\item If $a<p_2$, then $F'>0$ on $(p_2,m)$, where table $(*)$ allows only
$1\to2$ switches. If $H$ wins strictly at $p_2$, then $g(p_2)=0$; product~1
just right of $p_2$ would require $g>0$ with $g'=\delta g/b_1-F'<0$, which
is impossible starting from $g(p_2)=0$. So product~2 continues to the right
of $p_2$ up to $m$, and beyond $m$ only one $2\to1$ switch is possible. If
instead $L$ wins strictly at $p_2$, it wins on all of $[p_1,p_2]$ by (iii);
then at most one $1\to2$ switch occurs on $(p_2,m)$ and at most one $2\to1$
switch on $(m,1]$.
\item If $L(p_2)=H(p_2)$, then $F'(p_2)>0$, $L'(p_2)=0$, and
$H'(p_2)=F'(p_2)/(\delta+\lambda D(p_2))>0$, so $L$ wins just left of $p_2$,
while the $g'$ argument gives product~2 just right. This is a single
$1\to2$ switch.
\end{itemize}
In every case product~2 is used on at most one interval. At an attractive
peak $m\in(p_1,p_2)$, holding uses the mixture $(m-p_1)/\Delta p$.
\end{proof}

\begin{lemma}[Holding points]\label{lem:holding}
In problem \eqref{eq:1d-problem}, if $x^*\in(p_1,p_2)$ and holding at $x^*$
forever is optimal from $x^*$, then $F'(x^*)=0$ and $F''(x^*)\leq0$.
\end{lemma}

\begin{proof}
Holding yields $F(x^*)/\delta$. From $x^*$, both directions are available at
positive speed, and every state near $x^*$ can again be held. If
$F'(x^*)\neq0$, move for a short time $\tau$ in the direction in which $F$
increases and hold thereafter. The payoff changes by
$\int_0^\tau e^{-\delta t}[F(x(t))-F(x^*)]\,dt
+e^{-\delta\tau}[F(x(\tau))-F(x^*)]/\delta>0$, a contradiction. If
$F'(x^*)=0$ and $F''(x^*)>0$, then $x^*$ is a strict local minimum of $F$,
and the same deviation in either direction is profitable.
\end{proof}

\subsection{Proofs of the main single-interval statements}

For the following proofs, fix $y$ and use $D(x)=D_A(x,y)$ and
$F(x)=F_A(x,y)$.

\begin{proof}[Proof of Proposition~\ref{prop:frozen-reduction}]
Fix the rival's reputation $y$. The transformation
\eqref{eq:duopoly-transform} gives running payoff $F$ and drift
$\dot x=\lambda D(x)(p_1+u\Delta p-x)$. By
Lemma~\ref{lem:control}, it suffices to show that $F$ is either
nonincreasing, or decreases, then increases, then decreases (with absent
segments allowed), with $F''\neq0$ at interior zeros of $F'$. Here
$D\in(0,1)$ is smooth and $D'>0$ by
Lemma~\ref{lem:marginal-demand-representation}.

Let $r=D''/D'$, which is nonincreasing by assumption, and set
$\varphi(x)=(q-x)r(x)$. Differentiating $F$,
\begin{equation*}
F'(x)=\frac{\Delta c}{\Delta p}\Bigl[(q-x)D'(x)-D(x)-\frac\delta\lambda\Bigr],
\qquad
F''(x)=\frac{\Delta c}{\Delta p}D'(x)\bigl[\varphi(x)-2\bigr].
\end{equation*}
On $[q,1]$ (if nonempty), $F'<0$ because $D>0$ and $D'>0$. On
$[0,\min\{q,1\})$, the set where $r>0$ is an initial segment, because $r$ is
nonincreasing. On it, $\varphi$ is strictly decreasing: if $u<v$ and
$r(v)>0$, then $r(u)(q-u)\geq r(v)(q-u)>r(v)(q-v)$. Off it,
$\varphi\leq0<2$. Hence $\varphi-2$ changes sign at most once, from positive
to negative, and vanishes at most at one point $c$. Consequently, $F'$ is
strictly increasing on $[0,c]$, strictly decreasing on $[c,\min\{q,1\}]$,
and negative on $[q,1]$. The set $\{F'>0\}$ is therefore an interval $I$. If
$I$ is empty, $F$ is nonincreasing. If $I$ is nonempty, then $F'(c)>0$, so
neither endpoint of $I$ equals $c$, and $F''\neq0$ at both endpoints. This
is the required shape.
\end{proof}

\begin{proof}[Proof of Corollary~\ref{cor:single-seller}]
The proof of Proposition~\ref{prop:frozen-reduction} applies verbatim with
$D=S_{p_0}$, which satisfies $S_{p_0}'>0$ and is strictly log-concave by
Lemma~\ref{lem:single-seller-logconcave}.
\end{proof}

\begin{remark}[A weaker sufficient condition]
The proof only needs $F''$ to change sign at most once, from positive to
negative, on $[0,\min\{q,1\})$. Writing $\psi=(q-x)D''-2D'$, at any zero of
$\psi$ we have $D''=2D'/(q-x)>0$ and $\psi'=(q-x)D'''-3D''$, so $\psi'<0$
there exactly when $D'''<\tfrac32(D'')^2/D'$. Hence the condition
\begin{equation*}
D'D'''<\tfrac32\,(D'')^2\qquad\text{at points where }D''>0,
\end{equation*}
which says that $(D')^{-1/2}$ is convex where $D$ is convex, also suffices:
a $C^1$ function whose zeros all have negative derivative crosses zero at
most once, from above. Log-concavity is the stronger condition
$D'D'''\leq(D'')^2$.
\end{remark}

\section{Equilibrium proofs}\label{app:equilibrium}

\subsection{Stationary equilibrium}\label{app:stationary}

\begin{proof}[Proof of Proposition~\ref{prop:stationary-open-loop}]
Fix $u_B^*=(z-p_1)/\Delta p$. Starting from $y(0)=z$, the rival's dynamics
are $\dot y=\lambda D_B(x,y)(z-y)$. Uniqueness gives $y(t)=z$ under every
admissible deviation by $A$. Seller $A$ therefore faces exactly the
frozen-rival problem. Symmetry gives the identical problem for seller $B$.

By Proposition~\ref{prop:frozen-reduction}, the frozen transformed payoff
is nonincreasing or has the valley-peak shape used in
Lemma~\ref{lem:control}. The shape argument gives nondegenerate extrema
whenever the increasing segment is nonempty. An interior holding point
must consequently be the peak: a point with nonzero derivative admits a
profitable small move, a valley admits a move in either direction, and an
isolated horizontal point of a decreasing payoff admits a profitable move
to the left. Thus holding can be optimal only if $z=m(z)$, equivalently
\eqref{eq:stationary-peak-test}.

Suppose now that $z=m(z)$. Apply Step~2 of the proof of
Lemma~\ref{lem:control} with $h=m=z$ and the rival fixed at $z$. On
$[p_1,z]$, the transformed value is $\max\{L_z,H_z\}$, where $L_z$ is
the value of product~1 forever and $H_z$ is the value of product~2 until
reaching $z$, followed by holding. Excursions above $z$ cannot improve
the value because $F_A(x,z)\leq F_A(z,z)$ for every $x\geq z$.
At the initial state $z$, $H_z(z)=F_A(z,z)/\delta$. Hence
\begin{equation*}
W_A^{\mathrm{fr}}(z;z)
=\max\left\{L_z(z),\frac{F_A(z,z)}{\delta}\right\},
\end{equation*}
where $W_A^{\mathrm{fr}}$ denotes the optimal transformed frozen-rival
value. Holding is a global best response exactly when the inequality in
\eqref{eq:stationary-exact-test} holds. The payoff transformation adds
the same constant to every deviation payoff, and symmetry completes the
equilibrium argument.
\end{proof}

For the general holding test \eqref{eq:general-stationary-test}, the same
argument freezes the rival under every unilateral deviation. The holding
payoff is $F_A(x^*,y^*)/\delta$ for $A$ and $F_A(y^*,x^*)/\delta$ for $B$.
Equality with each optimal frozen-rival value is therefore necessary and
sufficient, including at boundary states. No shape restriction is needed
for this equivalence.

\subsection{Verification for a moving rival}\label{app:verification}

\begin{proof}[Verification of \eqref{eq:open-loop-verification}]
Along every admissible deviation against $u_B^*(t)$, the chain rule and
\eqref{eq:open-loop-verification} give
\begin{equation*}
\frac{d}{dt}\bigl[e^{-\delta t}U_A(t,x(t),y(t))\bigr]
\leq-e^{-\delta t}F_A(x(t),y(t)).
\end{equation*}
Integrating to $T$ and letting $T\to\infty$, boundedness of $U_A$ gives
an upper bound $U_A(0,x_0,y_0)$ for the transformed payoff. Equality holds
along a proposed path whose control attains the maximum almost everywhere.
The constant in \eqref{eq:duopoly-transform} preserves the ranking of
controls. Applying the same argument to $B$ proves mutual global best
responses and hence open-loop Nash equilibrium from the initial state.
\end{proof}

\section{Feedback comparison}\label{app:feedback}

This appendix records the feedback formulation for comparison with the
open-loop problem studied in the main text.

Given $B$'s Markov strategy $u_B(x,y)$, seller $A$'s value satisfies, over
one period,
\begin{align*}
V_A(x,y)={}&D_A\max_{u}\Bigl\{(R-c_1-u\Delta c)dt
+\gamma^{dt}V_A\bigl(x+\lambda(p_1+u\Delta p-x)dt,\ y\bigr)\Bigr\}\\
&+D_B\,\gamma^{dt}V_A\bigl(x,\ y+\lambda(p_1+u_B\Delta p-y)dt\bigr)+o(dt).
\end{align*}
A first-order Taylor expansion and $D_A+D_B=1$ give, after dividing by $dt$
and letting $dt\to0$, the Hamilton--Jacobi--Bellman (HJB) equation
\begin{equation}\label{eq:duopoly-HJB}
\delta V_A=D_A\max_{u\in[0,1]}\Bigl\{R-c_1-u\Delta c
+\lambda(p_1+u\Delta p-x)\,\partial_xV_A\Bigr\}
+\lambda D_B\bigl(p_1+u_B(x,y)\Delta p-y\bigr)\partial_yV_A,
\end{equation}
together with the symmetric equation for $V_B$.

\paragraph{Best response.}
The expression inside the maximum is linear in $u$ with slope
$D_A[-\Delta c+\lambda\Delta p\,\partial_xV_A]$. Wherever $V_A$ is
differentiable, product~2 is therefore strictly preferred if and only if
\begin{equation*}
\partial_xV_A(x,y)>\frac{\Delta c}{\lambda\Delta p}.
\end{equation*}
Away from indifference, optimal controls are bang-bang. At indifference,
mixtures may be optimal; they need not hold the state fixed. A holding
mixture is feasible on $[p_1,p_2]$, where $u=(x-p_1)/\Delta p$ keeps $x$
fixed, but its optimality requires a separate argument.

\begin{definition}[Feedback Nash equilibrium]
A pair of Markov strategies $u_A^*(x,y)$, $u_B^*(x,y)$, measurable from
$[0,1]^2$ to $[0,1]$, is a feedback Nash equilibrium if, at every state,
$u_A^*$ attains the maximum in \eqref{eq:duopoly-HJB} with $u_B=u_B^*$,
where $V_A$ is seller $A$'s optimal value against $u_B^*$, and symmetrically
for $B$. Where $V_A$ is not differentiable, \eqref{eq:duopoly-HJB} is
understood in the viscosity sense.
\end{definition}

Two technical points arise. First, bang-bang feedback strategies make the
drift discontinuous, so trajectories must be defined in the Filippov sense,
or strategies restricted to finitely many switching curves on which sliding
is resolved by holding mixtures. Second, existence of feedback equilibria in
differential games is not available in general. An alternative route is to
take fluid limits of Markov perfect equilibria of the discrete game.

\paragraph{The strategic term.}
Feedback and open-loop equilibria differ in one term. At points where $V_A$
is twice differentiable and $u_B^*$ is differentiable in $x$,
differentiating \eqref{eq:duopoly-HJB} with respect to $x$ produces
\begin{equation*}
\lambda D_B(x,y)\,\Delta p\,\partial_xu_B^*(x,y)\,\partial_yV_A(x,y),
\end{equation*}
which has no counterpart in the open-loop costate equation for $\xi$,
where $u_B^*$ depends on time only. Since $D_A$ is nonincreasing in $y$, one
expects $\partial_yV_A\leq0$. If the rival raises quality when $A$'s
reputation improves ($\partial_xu_B^*>0$), this term lowers $A$'s marginal
value of reputation and discourages product~2; if the rival backs off, it
encourages product~2. With bang-bang strategies, $\partial_xu_B^*$ is
concentrated on the rival's switching curve.

\begin{thebibliography}{9}
\bibitem{badia2023}
F.~G.~Bad\'ia, H.~Lee, and C.~Sang\"uesa (2023).
Log concavity preservation by beta operator based on probability tools.
\emph{Journal of Mathematical Analysis and Applications} 525(1), 127122.
\href{https://doi.org/10.1016/j.jmaa.2023.127122}{doi:10.1016/j.jmaa.2023.127122}.

\bibitem{hu2026}
X.~Hu, M.~Jagadeesan, M.~I.~Jordan, and J.~Steinhardt (2026).
Incentivizing high-quality content in online recommender systems.
In \emph{7th Symposium on Foundations of Responsible Computing (FORC
2026)}, LIPIcs 368, 15:1--15:18.
\end{thebibliography}

\end{document}
